\section{多时钟设计怎样定义时序关系}

一块真实芯片可能同时存在 CPU 核心时钟、总线时钟、UART 参考时钟和外部设备
时钟。静态时序分析必须知道每个时钟从哪里来、周期多长、彼此是否有固定相位
关系。没有声明关系的路径不会自动变安全，只会让分析工具缺少判断依据。

\topic{生成时钟不能只写一个频率名字}
PLL、分频器和时钟选择器会从源时钟产生新时钟。生成时钟需要记录源、倍频、
分频、相位和传播路径。若两个时钟来自不同振荡器，应按异步时钟组处理；若它们
来自同一 PLL，仍要依据实际相位关系分析跨域路径。

\begin{osgridlongtable}{L{0.20\linewidth}L{0.30\linewidth}L{0.40\linewidth}}
关系 & 例子 & 分析方式 \\ \hline
\endfirsthead
关系 & 例子 & 分析方式 \\ \hline
\endhead
同一时钟 & 两个寄存器都由 CLK\_A 采样 & 直接检查建立与保持 \\ \hline
同步生成 & CLK\_B 是 CLK\_A 的二分频 & 建立生成时钟和确定相位关系 \\ \hline
异步 & 两个独立晶振 & 使用 CDC 结构，不把路径当普通同步路径 \\ \hline
条件互斥 & 复用器任一时刻只选一路时钟 & 证明互斥条件并使用专用时钟切换结构 \\ \hline
\end{osgridlongtable}

\section{时序例外为什么必须带结构证明}

\term{false path}{功能上永远不会被采样的路径}和
\term{multicycle path}{允许跨多个周期到达的路径}会改变工具的检查要求。
错误例外可能让真正违例从报告中消失，因此不能为了让报表变绿而添加。

\topic{多周期路径同时影响建立与保持}
若一条路径允许两个周期完成，建立检查可以移动到第二个采样边沿；保持检查也要
按照工具语义同步调整。只放宽建立而不理解保持边沿，可能制造新的最短路径
违例。设计评审需要同时展示源寄存器、目标寄存器、使能条件和实际采样周期。

\begin{failurebox}
“仿真没有失败”不能证明时序例外正确。零延迟 RTL 仿真通常看不到布线延迟、
时钟偏斜和 PVT 角落。例外需要由电路结构、协议状态和形式检查共同支持。
\end{failurebox}

\section{复杂 CDC 怎样传递多位状态}

单 bit 慢速电平可以使用两级同步器。连续数据流、突发命令和多位计数器需要
协议，否则接收端可能看到混合字或丢失事件。

\topic{握手让数据在跨域期间保持不变}
源域先把数据写入保持寄存器，再翻转 request；目标域同步 request、采样稳定
数据并返回 acknowledge；源域收到 acknowledge 后才允许修改数据。吞吐低于
异步 FIFO，但状态容易逐步检查。

\topic{异步 FIFO 用 Gray 指针减少同时变化的位}
二进制指针从 \(\texttt{0111}\) 加一到 \(\texttt{1000}\) 时四位同时变化；
跨域独立采样可能得到混合值。Gray 编码让相邻计数只改变一位，再把同步后的
读写指针用于空满判断。存储阵列、指针同步和空满逻辑仍要分别验证。

\begin{osgridlongtable}{L{0.21\linewidth}L{0.31\linewidth}L{0.38\linewidth}}
载荷 & 常用结构 & 仍需检查 \\ \hline
\endfirsthead
载荷 & 常用结构 & 仍需检查 \\ \hline
\endhead
慢速单 bit 电平 & 两级同步器 & 延迟、MTBF、复位值 \\ \hline
单次事件 & toggle 或 request/ack & 不丢事件、不重复计数 \\ \hline
稳定多 bit 配置 & 数据保持加握手 & 保持窗口与确认返回 \\ \hline
连续数据流 & 异步 FIFO & Gray 指针、空满边界、双端复位 \\ \hline
\end{osgridlongtable}

\section{复位域与低功耗状态怎样交叉}

每个时钟域的复位释放需要与本域时钟对齐。电源域关闭以后，输出可能失去有效
电平；跨域接口还需要隔离单元、保持寄存器和确定的上电顺序。RDC 检查关注的
是复位关系，不能被普通 CDC 报告替代。

\topic{验证报告应回到具体路径}
检查复杂时序时，报告至少要能定位源寄存器、目标寄存器、时钟、约束、跨域
结构和复位关系。每个被豁免的路径都应说明为什么不会在硅片上被错误采样。
这些工作属于芯片和 FPGA 实现阶段，当前 QEMU 教学 OS 不会直接执行它们。
